\section{Results}
\label{sec:results}

\subsection{Dataset and protocol overview}

The common benchmark included four outcome tasks from three sources. Vigitel obesity and current smoking share the Vigitel source and are not independent cohorts. The analytic samples contained 851 of 5,000 positive obesity records, 83 deaths among 334 COVID-19 admissions, 616 of 5,000 current smokers, and 267 deaths among 5,000 kidney-cohort patients. Every method used the same held-out real folds within a task. The augmentation export contains 1,620 fold--classifier rows (4 tasks, 9 methods, 3 classifiers, 3 repetitions, and 5 folds); no result was aggregated before export.

\subsection{Real-only and conventional imbalance baselines}

Table~\ref{tab:common-benchmark} reports the full augmentation comparison. Real-only training had the highest mean AUPRC for Vigitel obesity (0.287), Vigitel smoking (0.300), and kidney mortality (0.122). Random oversampling had the highest COVID-19 AUPRC (0.355). Threshold-dependent conclusions differed: for kidney mortality, real-only training had mean positive-class recall below 0.001 at the fixed 0.5 threshold, whereas random undersampling increased recall to 0.617 but increased the Brier score from 0.050 to 0.225. These results show why ranking, classification at a fixed threshold, and calibration are reported separately.

\begin{longtable}{p{.17\textwidth}p{.24\textwidth}rrrrr}
\caption{Common augmentation benchmark. Values are means over 3 repeated 5-fold partitions and 3 classifiers (45 paired fold--classifier observations per method and task). Held-out real folds retain natural prevalence; uncertainty and paired inference are reported separately.}\label{tab:common-benchmark}\\
\toprule
Task & Method & AUPRC & Macro-F1 & Recall & Brier & $n$ \\
\midrule
\endfirsthead
\toprule
Task & Method & AUPRC & Macro-F1 & Recall & Brier & $n$ \\
\midrule
\endhead
COVID-19 & Class weighting & 0.348 & 0.530 & 0.373 & 0.234 & 45 \\
COVID-19 & Random oversampling & 0.355 & 0.538 & 0.374 & 0.231 & 45 \\
COVID-19 & Random undersampling & 0.332 & 0.500 & 0.512 & 0.279 & 45 \\
COVID-19 & Real only & 0.351 & 0.493 & 0.113 & 0.198 & 45 \\
COVID-19 & SMOTE & 0.326 & 0.507 & 0.204 & 0.218 & 45 \\
COVID-19 & SepAware CTGAN & 0.323 & 0.534 & 0.362 & 0.239 & 45 \\
COVID-19 & SepAware TVAE & 0.343 & 0.508 & 0.178 & 0.215 & 45 \\
COVID-19 & Standard CTGAN & 0.311 & 0.521 & 0.355 & 0.240 & 45 \\
COVID-19 & Standard TVAE & 0.350 & 0.510 & 0.183 & 0.214 & 45 \\
Kidney mortality & Class weighting & 0.106 & 0.523 & 0.371 & 0.156 & 45 \\
Kidney mortality & Random oversampling & 0.103 & 0.528 & 0.367 & 0.155 & 45 \\
Kidney mortality & Random undersampling & 0.117 & 0.487 & 0.617 & 0.225 & 45 \\
Kidney mortality & Real only & 0.122 & 0.487 & 0.000 & 0.050 & 45 \\
Kidney mortality & SMOTE & 0.096 & 0.520 & 0.256 & 0.131 & 45 \\
Kidney mortality & SepAware CTGAN & 0.110 & 0.526 & 0.175 & 0.099 & 45 \\
Kidney mortality & SepAware TVAE & 0.103 & 0.508 & 0.054 & 0.065 & 45 \\
Kidney mortality & Standard CTGAN & 0.104 & 0.522 & 0.140 & 0.094 & 45 \\
Kidney mortality & Standard TVAE & 0.103 & 0.503 & 0.042 & 0.062 & 45 \\
Vigitel & Class weighting & 0.284 & 0.572 & 0.476 & 0.201 & 45 \\
Vigitel & Random oversampling & 0.281 & 0.573 & 0.486 & 0.202 & 45 \\
Vigitel & Random undersampling & 0.277 & 0.551 & 0.618 & 0.230 & 45 \\
Vigitel & Real only & 0.287 & 0.462 & 0.009 & 0.134 & 45 \\
Vigitel & SMOTE & 0.233 & 0.503 & 0.204 & 0.176 & 45 \\
Vigitel & SepAware CTGAN & 0.256 & 0.527 & 0.254 & 0.175 & 45 \\
Vigitel & SepAware TVAE & 0.252 & 0.492 & 0.089 & 0.153 & 45 \\
Vigitel & Standard CTGAN & 0.247 & 0.511 & 0.219 & 0.176 & 45 \\
Vigitel & Standard TVAE & 0.253 & 0.494 & 0.096 & 0.153 & 45 \\
Vigitel smoking & Class weighting & 0.295 & 0.574 & 0.452 & 0.165 & 45 \\
Vigitel smoking & Random oversampling & 0.291 & 0.582 & 0.467 & 0.165 & 45 \\
Vigitel smoking & Random undersampling & 0.284 & 0.559 & 0.654 & 0.210 & 45 \\
Vigitel smoking & Real only & 0.300 & 0.500 & 0.036 & 0.099 & 45 \\
Vigitel smoking & SMOTE & 0.229 & 0.541 & 0.192 & 0.134 & 45 \\
Vigitel smoking & SepAware CTGAN & 0.249 & 0.565 & 0.288 & 0.143 & 45 \\
Vigitel smoking & SepAware TVAE & 0.255 & 0.526 & 0.099 & 0.111 & 45 \\
Vigitel smoking & Standard CTGAN & 0.217 & 0.536 & 0.216 & 0.145 & 45 \\
Vigitel smoking & Standard TVAE & 0.253 & 0.523 & 0.095 & 0.111 & 45 \\
\bottomrule
\end{longtable}


\begin{figure}[htbp]
\centering
\includegraphics[width=.96\textwidth]{common_benchmark_auprc.pdf}
\caption{Held-out-real AUPRC in the augmentation regime. Bars are means over 3 repeated 5-fold partitions and 3 classifiers. Methods within each task use identical test folds; corrected paired comparisons are shown in Figure~\ref{fig:primary-contrasts}.}
\label{fig:common-auprc}
\end{figure}

\subsection{Standard generators and \sepa{} selection}

For CTGAN, mean AUPRC differences between \sepa{} and standard sampling were 0.012 for COVID-19, 0.009 for Vigitel obesity, 0.031 for Vigitel smoking, and 0.005 for kidney mortality. The Nadeau--Bengio corrected intervals were $[-0.055,0.079]$, $[-0.004,0.021]$, $[0.015,0.048]$, and $[-0.007,0.017]$, respectively. Only smoking remained significant after Holm adjustment ($p_{\mathrm{adj}}=.009$). TVAE differences ranged from $-0.007$ to 0.003, and all corrected intervals included zero. The inferential unit was the mean difference across classifiers within each repeated-CV fold ($n=15$ fold summaries per contrast), not 45 fold--classifier observations.

\begin{figure}[htbp]
\centering
\includegraphics[width=.82\textwidth]{primary_contrasts.pdf}
\caption{Primary paired AUPRC contrasts in real-plus-synthetic augmentation. Points are mean \sepa{} minus standard-generator differences; intervals use the Nadeau--Bengio repeated-cross-validation correction. Eight task--generator tests were adjusted together by Holm's method.}
\label{fig:primary-contrasts}
\end{figure}

The selected CTGAN condition did not surpass the strongest conventional method on any task. Its mean AUPRC was 0.323 for COVID-19, 0.256 for Vigitel obesity, 0.249 for Vigitel smoking, and 0.110 for kidney mortality. The corresponding best conventional values were 0.355, 0.287, 0.300, and 0.122. SepAware therefore improved one generator comparison reliably without establishing superiority over simpler training strategies.

\subsection{TSTR utility}
\label{sec:tstr-results}

The common TSTR experiment trained each classifier only on a synthetic dataset whose class counts matched the corresponding real training fold and evaluated on the untouched real test fold. Standard and selected variants shared candidate pools. Table~\ref{tab:tstr-summary} reports all tasks and both generators; these substitution results are kept separate from augmentation.

\begin{longtable}{p{.20\textwidth}p{.24\textwidth}rrrr}
\caption{TSTR performance. Values are means over 3 repeated 5-fold partitions and 3 classifiers (45 paired observations per condition and task). Synthetic class counts match each real training fold; evaluation uses held-out real records.}\label{tab:tstr-summary}\\
\toprule
Task & Synthetic training data & AUPRC & Macro-F1 & Recall & Brier \\
\midrule
\endfirsthead
\toprule
Task & Synthetic training data & AUPRC & Macro-F1 & Recall & Brier \\
\midrule
\endhead
COVID-19 hospital mortality & SepAware CTGAN & 0.333 & 0.517 & 0.194 & 0.229 \\
COVID-19 hospital mortality & SepAware TVAE & 0.324 & 0.478 & 0.498 & 0.385 \\
COVID-19 hospital mortality & Standard CTGAN & 0.279 & 0.459 & 0.068 & 0.214 \\
COVID-19 hospital mortality & Standard TVAE & 0.332 & 0.491 & 0.468 & 0.347 \\
Kidney-cohort 365-day mortality & SepAware CTGAN & 0.080 & 0.491 & 0.007 & 0.059 \\
Kidney-cohort 365-day mortality & SepAware TVAE & 0.100 & 0.519 & 0.368 & 0.176 \\
Kidney-cohort 365-day mortality & Standard CTGAN & 0.063 & 0.486 & 0.000 & 0.052 \\
Kidney-cohort 365-day mortality & Standard TVAE & 0.102 & 0.524 & 0.336 & 0.160 \\
Vigitel current smoking & SepAware CTGAN & 0.279 & 0.601 & 0.319 & 0.128 \\
Vigitel current smoking & SepAware TVAE & 0.179 & 0.443 & 0.689 & 0.419 \\
Vigitel current smoking & Standard CTGAN & 0.139 & 0.467 & 0.000 & 0.109 \\
Vigitel current smoking & Standard TVAE & 0.183 & 0.428 & 0.700 & 0.432 \\
Vigitel obesity & SepAware CTGAN & 0.279 & 0.548 & 0.175 & 0.148 \\
Vigitel obesity & SepAware TVAE & 0.200 & 0.348 & 0.847 & 0.572 \\
Vigitel obesity & Standard CTGAN & 0.191 & 0.454 & 0.001 & 0.143 \\
Vigitel obesity & Standard TVAE & 0.223 & 0.502 & 0.447 & 0.308 \\
\bottomrule
\end{longtable}


For CTGAN, selection increased mean TSTR AUPRC from 0.279 to 0.333 for COVID-19, 0.191 to 0.279 for Vigitel obesity, 0.139 to 0.279 for smoking, and 0.063 to 0.080 for kidney mortality. TVAE changes were smaller and mixed: 0.332 to 0.324, 0.223 to 0.200, 0.183 to 0.179, and 0.102 to 0.100, respectively. These are descriptive substitution results and do not replace the prespecified augmentation contrasts.

\subsection{Calibration, coverage, diversity, fidelity, and privacy diagnostics}

In augmentation, CTGAN selection reduced positive-class coverage distance for every task: from 4.410 to 4.115 for COVID-19, 1.977 to 1.723 for Vigitel obesity, 3.833 to 3.586 for smoking, and 1.744 to 1.595 for kidney mortality. It also reduced within-positive-class diversity from 11.332 to 7.888, 5.478 to 4.356, 7.340 to 6.310, and 6.827 to 4.425, respectively. Thus closer coverage was consistently accompanied by a more concentrated selected sample.

In augmentation, exact-match rates were zero or below 0.0005, and membership-proxy AUROCs ranged from 0.499 to 0.508 for standard and selected CTGAN. In TSTR, kidney TVAE had much higher exact-match rates: 0.105 for standard TVAE and 0.134 after selection. TSTR detection AUROC was also high for several conditions, reaching 0.995--0.998 for kidney CTGAN and at least 0.979 for kidney TVAE. These results show readily detectable distributional differences and, for kidney TVAE, material duplication risk. Membership-proxy AUROCs remained near 0.5, but none of these measures establishes anonymity. The full marginal similarity, correlation similarity, detection, coverage, diversity, exact-match, and membership-proxy matrix is provided as a machine-readable supplementary table.

\subsection{Dataset-specific anchor and factorial analyses}

Clinical-anchor analyses were not part of the common benchmark. The COVID-19 DACPF experiment required age, cancer, and respiratory-admission variables; comparable anchors were not prespecified for smoking or kidney mortality. Applying invented substitutes after seeing the data would have introduced outcome-guided tuning.

In the five-fold COVID-19 DACPF augmentation run, disease-anchor filtering produced mean Macro-F1 of 0.577 for CatBoost and 0.534 for XGBoost, compared with 0.526 for each classifier under standard CTGAN augmentation. Age-only and combined filters did not improve on standard augmentation. The result is dataset- and classifier-specific.

The blocked $2^2$ anchor-by-separability TSTR analysis was likewise limited to Vigitel obesity and COVID-19. The fold-paired Vigitel separability effect was 0.102 Macro-F1 (95\% CI 0.077--0.128; exploratory $p<.001$). The COVID-19 estimate was $-0.018$ (95\% CI $-0.063$--0.027; $p=.332$). Anchor and interaction intervals included zero for both tasks. Overlapping training folds and shared candidate pools limit these exploratory intervals; they are not independent replications.

\subsection{Controlled simulation}

The simulation varied prevalence and class overlap independently and is not counted as a health-data task. SepAware-SMOTE exceeded ordinary SMOTE in each of nine prevalence--overlap cells for logistic regression, but mean differences were small (0.001--0.024). In pooled heteroscedasticity-robust regression, neither the method contrast nor its overlap and prevalence interactions was distinguishable from zero (all $p>.91$). The simulation therefore supports separating prevalence from overlap but does not locate a stable region in which selection improves predictive performance.

\subsection{Sensitivity and statistical interpretation}

Classifier-by-repetition effects are exported with the primary statistics. Their variation is consistent with the wide corrected intervals for seven of eight primary contrasts. Results from Vigitel obesity and smoking cannot be treated as two independent source replications. The fixed-$\lambda$ grid is retained only as an exploratory development record because configurations were compared on the same folds used for reporting.
